This dissertation investigates the theoretical and statistical foundations of graviton-mass and polarization-mode inference in pulsar timing arrays (PTAs), elucidating the impact of likelihood approximations, frequency compression, and dispersive response templates. Based on linear massive gravitational wave propagation and the Fierz--Pauli field theory, we develop a pioneering paired diagnostic framework comparing complex Fourier coefficients, quadratic estimators, and exact Gaussian controls with matching physical moments.

The results reveal that the standard normal approximation applied to physical estimators fails simulation-based calibration (SBC), exhibiting persistent rejections under Holm-Bonferroni correction ($p \le 0.005$) and severe tail undercoverage. The proposed Gaussian control isolates and quantifies for the first time the net information penalty of spectral compression, removing on average $0.06211 \pm 0.00477$ and $0.10631 \pm 0.00695$ nat of information about the mass across two array configurations based on the MeerKAT catalog. Furthermore, reference-frequency approximations produce small mean shifts ($< 0.005$) in the $Q_{0.95}(u)$ quantile primarily because prior compression severely degrades the sensitivity to spectral features.

In addition, we prove analytically that under a uniform prior on $u = f_g T \in [0,1]$, the bound $m_g c^2 \le 0.95 h_{\mathrm P}/T$ emerges structurally even under a constant likelihood with zero Kullback--Leibler information gain ($D_{\rm KL} = 0$), establishing that learning metrics must accompany reported mass limits. In the scalar sector, an exact angular degeneracy $\Gamma_0 = \Gamma_T/2$ is proven at the propagation threshold. This work establishes essential diagnostic criteria and practical guidelines for future observational analyses in international PTA consortia.
